% Options for packages loaded elsewhere
\PassOptionsToPackage{unicode}{hyperref}
\PassOptionsToPackage{hyphens}{url}
\PassOptionsToPackage{dvipsnames,svgnames,x11names}{xcolor}
\PassOptionsToPackage{space}{xeCJK}
\documentclass[
]{article}
\usepackage{xcolor}
\usepackage[margin=2.2cm]{geometry}
\usepackage{amsmath,amssymb}
\setcounter{secnumdepth}{-\maxdimen} % remove section numbering
\usepackage{iftex}
\ifPDFTeX
  \usepackage[T1]{fontenc}
  \usepackage[utf8]{inputenc}
  \usepackage{textcomp} % provide euro and other symbols
\else % if luatex or xetex
  \usepackage{unicode-math} % this also loads fontspec
  \defaultfontfeatures{Scale=MatchLowercase}
  \defaultfontfeatures[\rmfamily]{Ligatures=TeX,Scale=1}
\fi
\usepackage{lmodern}
\ifPDFTeX\else
  % xetex/luatex font selection
  \setmainfont[]{Microsoft YaHei}
  \setmonofont[]{Consolas}
  \ifXeTeX
    \usepackage{xeCJK}
    \setCJKmainfont[]{Microsoft YaHei}
  \fi
  \ifLuaTeX
    \usepackage[]{luatexja-fontspec}
    \setmainjfont[]{Microsoft YaHei}
  \fi
\fi
% Use upquote if available, for straight quotes in verbatim environments
\IfFileExists{upquote.sty}{\usepackage{upquote}}{}
\IfFileExists{microtype.sty}{% use microtype if available
  \usepackage[]{microtype}
  \UseMicrotypeSet[protrusion]{basicmath} % disable protrusion for tt fonts
}{}
\makeatletter
\@ifundefined{KOMAClassName}{% if non-KOMA class
  \IfFileExists{parskip.sty}{%
    \usepackage{parskip}
  }{% else
    \setlength{\parindent}{0pt}
    \setlength{\parskip}{6pt plus 2pt minus 1pt}}
}{% if KOMA class
  \KOMAoptions{parskip=half}}
\makeatother
\usepackage{color}
\usepackage{fancyvrb}
\newcommand{\VerbBar}{|}
\newcommand{\VERB}{\Verb[commandchars=\\\{\}]}
\DefineVerbatimEnvironment{Highlighting}{Verbatim}{commandchars=\\\{\}}
% Add ',fontsize=\small' for more characters per line
\newenvironment{Shaded}{}{}
\newcommand{\AlertTok}[1]{\textcolor[rgb]{1.00,0.00,0.00}{\textbf{#1}}}
\newcommand{\AnnotationTok}[1]{\textcolor[rgb]{0.38,0.63,0.69}{\textbf{\textit{#1}}}}
\newcommand{\AttributeTok}[1]{\textcolor[rgb]{0.49,0.56,0.16}{#1}}
\newcommand{\BaseNTok}[1]{\textcolor[rgb]{0.25,0.63,0.44}{#1}}
\newcommand{\BuiltInTok}[1]{\textcolor[rgb]{0.00,0.50,0.00}{#1}}
\newcommand{\CharTok}[1]{\textcolor[rgb]{0.25,0.44,0.63}{#1}}
\newcommand{\CommentTok}[1]{\textcolor[rgb]{0.38,0.63,0.69}{\textit{#1}}}
\newcommand{\CommentVarTok}[1]{\textcolor[rgb]{0.38,0.63,0.69}{\textbf{\textit{#1}}}}
\newcommand{\ConstantTok}[1]{\textcolor[rgb]{0.53,0.00,0.00}{#1}}
\newcommand{\ControlFlowTok}[1]{\textcolor[rgb]{0.00,0.44,0.13}{\textbf{#1}}}
\newcommand{\DataTypeTok}[1]{\textcolor[rgb]{0.56,0.13,0.00}{#1}}
\newcommand{\DecValTok}[1]{\textcolor[rgb]{0.25,0.63,0.44}{#1}}
\newcommand{\DocumentationTok}[1]{\textcolor[rgb]{0.73,0.13,0.13}{\textit{#1}}}
\newcommand{\ErrorTok}[1]{\textcolor[rgb]{1.00,0.00,0.00}{\textbf{#1}}}
\newcommand{\ExtensionTok}[1]{#1}
\newcommand{\FloatTok}[1]{\textcolor[rgb]{0.25,0.63,0.44}{#1}}
\newcommand{\FunctionTok}[1]{\textcolor[rgb]{0.02,0.16,0.49}{#1}}
\newcommand{\ImportTok}[1]{\textcolor[rgb]{0.00,0.50,0.00}{\textbf{#1}}}
\newcommand{\InformationTok}[1]{\textcolor[rgb]{0.38,0.63,0.69}{\textbf{\textit{#1}}}}
\newcommand{\KeywordTok}[1]{\textcolor[rgb]{0.00,0.44,0.13}{\textbf{#1}}}
\newcommand{\NormalTok}[1]{#1}
\newcommand{\OperatorTok}[1]{\textcolor[rgb]{0.40,0.40,0.40}{#1}}
\newcommand{\OtherTok}[1]{\textcolor[rgb]{0.00,0.44,0.13}{#1}}
\newcommand{\PreprocessorTok}[1]{\textcolor[rgb]{0.74,0.48,0.00}{#1}}
\newcommand{\RegionMarkerTok}[1]{#1}
\newcommand{\SpecialCharTok}[1]{\textcolor[rgb]{0.25,0.44,0.63}{#1}}
\newcommand{\SpecialStringTok}[1]{\textcolor[rgb]{0.73,0.40,0.53}{#1}}
\newcommand{\StringTok}[1]{\textcolor[rgb]{0.25,0.44,0.63}{#1}}
\newcommand{\VariableTok}[1]{\textcolor[rgb]{0.10,0.09,0.49}{#1}}
\newcommand{\VerbatimStringTok}[1]{\textcolor[rgb]{0.25,0.44,0.63}{#1}}
\newcommand{\WarningTok}[1]{\textcolor[rgb]{0.38,0.63,0.69}{\textbf{\textit{#1}}}}
\usepackage{longtable,booktabs,array}
\newcounter{none} % for unnumbered tables
\usepackage{calc} % for calculating minipage widths
% Correct order of tables after \paragraph or \subparagraph
\usepackage{etoolbox}
\makeatletter
\patchcmd\longtable{\par}{\if@noskipsec\mbox{}\fi\par}{}{}
\makeatother
% Allow footnotes in longtable head/foot
\IfFileExists{footnotehyper.sty}{\usepackage{footnotehyper}}{\usepackage{footnote}}
\makesavenoteenv{longtable}
\usepackage{graphicx}
\makeatletter
\newsavebox\pandoc@box
\newcommand*\pandocbounded[1]{% scales image to fit in text height/width
  \sbox\pandoc@box{#1}%
  \Gscale@div\@tempa{\textheight}{\dimexpr\ht\pandoc@box+\dp\pandoc@box\relax}%
  \Gscale@div\@tempb{\linewidth}{\wd\pandoc@box}%
  \ifdim\@tempb\p@<\@tempa\p@\let\@tempa\@tempb\fi% select the smaller of both
  \ifdim\@tempa\p@<\p@\scalebox{\@tempa}{\usebox\pandoc@box}%
  \else\usebox{\pandoc@box}%
  \fi%
}
% Set default figure placement to htbp
\def\fps@figure{htbp}
\makeatother
\setlength{\emergencystretch}{3em} % prevent overfull lines
\providecommand{\tightlist}{%
  \setlength{\itemsep}{0pt}\setlength{\parskip}{0pt}}
% texbook.py 注入的页眉/页脚样式（生成物，勿手改）
\usepackage{fancyhdr}
\pagestyle{fancy}
\fancyhf{}
\fancyhead[L]{\small\color{gray}Python 科学计算（新版教程）}
\fancyhead[R]{\small\color{gray}\thepage}
\renewcommand{\headrulewidth}{0.4pt}

% ===== 全局排版：行距/表格/图片/断行（M10 排版精修）=====
\usepackage{setspace}
\setstretch{1.15}
\renewcommand{\arraystretch}{1.18}
\setlength{\tabcolsep}{4pt}
\emergencystretch=1.5em
\clubpenalty=4000
\widowpenalty=4000
\RequirePackage{graphicx}
\setkeys{Gin}{width=\textwidth,keepaspectratio}

% 长表格用 \small 字体，减少列溢出
\usepackage{longtable}
\usepackage{etoolbox}
\AtBeginEnvironment{longtable}{\small}
% ===== 封面 / 章扉页（build/cover.tex；由 build 脚本自动引入，可直接改本文件定制）=====
% 不要在此删除 \maketitle：单章/附录 PDF 仍需要默认标题页；全书由 no_maketitle.tex 关闭。
\definecolor{courseblue}{RGB}{0, 82, 155}
\definecolor{coursegray}{RGB}{120, 120, 120}

\newcommand{\bookcover}[4]{%
  \begin{titlepage}
  \thispagestyle{empty}
  \centering
  \vspace*{1.6cm}
  {\color{courseblue}\rule{\linewidth}{4pt}}\\[16pt]
  {\Huge\bfseries #1}\\[10pt]
  {\Large\itshape #2}\\[14pt]
  {\color{courseblue}\rule{\linewidth}{1.2pt}}\\[22pt]
  {\large #3}\\[6pt]
  {\large #4}\\[10pt]
  {\small\color{coursegray}（Markdown 单一事实源 · Pandoc + XeLaTeX 自动排版）}\\[8pt]
  \vfill
  {\small\color{coursegray}配套：每章合订 PDF / 可编译 LaTeX 工程 / GitHub Pages 在线版}
  \end{titlepage}
}

\newcommand{\chaptercover}[2]{%
  \clearpage
  \thispagestyle{empty}
  \centering
  \vspace*{2.2cm}
  {\color{courseblue}\rule{\linewidth}{3pt}}\\[18pt]
  {\Huge\bfseries 第 #1 章}\\[18pt]
  {\LARGE\bfseries #2}\\[24pt]
  {\color{coursegray}\rule{0.55\linewidth}{0.5pt}}\\[16pt]
  {\small 本章配套：正文小节 / 上机 lab / 练习 quiz / 常见误区与技巧}
  \newpage
}
\usepackage{bookmark}
\IfFileExists{xurl.sty}{\usepackage{xurl}}{} % add URL line breaks if available
\urlstyle{same}
% fallback for those not using the hyperref driver hyperxmp:
\makeatletter
\@ifundefined{xmpquote}{\newcommand{\xmpquote}[1]{#1}}{}
\makeatother
\hypersetup{
  colorlinks=true,
  linkcolor={blue},
  filecolor={Maroon},
  citecolor={Blue},
  urlcolor={Blue},
  pdfcreator={LaTeX via pandoc}}

\author{}
\date{}

\begin{document}

\section{01 AI 编程与 Agent
基础}\label{ai-ux7f16ux7a0bux4e0e-agent-ux57faux7840}

\begin{quote}
本节回答三个问题：Agent 编程与''AI 帮我写代码''到底差在哪？一个 Agent
内部做了什么？哪些事必须自己管？ 阅读约 25 分钟；配套图 4
张，可直接运行的代码 1 段。
\end{quote}

\subsection{本节目标}\label{ux672cux8282ux76eeux6807}

\begin{enumerate}
\def\labelenumi{\arabic{enumi}.}
\tightlist
\item
  说清\textbf{代码补全、对话式编程、Agent
  编程}三者的差别，并解释''闭环''为什么是关键。
\item
  用\textbf{五件套}描述一个 agent
  harness（智能体框架）：模型、工具、上下文、权限、会话。
\item
  复述\textbf{一次 Agent 循环}的五个动作，并能把一个小任务拆成循环。
\item
  解释\textbf{上下文窗口与压缩}，知道怎么''省上下文''。
\item
  说出\textbf{权限三档}与审批机制，记住数据红线在哪里。
\item
  识别\textbf{四种典型失败模式}，并给出至少一条对策。
\end{enumerate}

\subsection{先修}\label{ux5148ux4fee}

\begin{itemize}
\tightlist
\item
  完成第 0 章的环境安装：会用终端（PowerShell / Bash）、会运行 python
  脚本；
\item
  学完第 1 章 NumPy：本节的演示案例用到 axis 与广播；
\item
  不需要任何机器学习或大模型背景。
\end{itemize}

\subsection{术语速查（中英对照）}\label{ux672fux8bedux901fux67e5ux4e2dux82f1ux5bf9ux7167}

后续章节会反复用到这些词，先混个脸熟即可：

{\def\LTcaptype{none} % do not increment counter
\begin{longtable}[]{@{}
  >{\raggedright\arraybackslash}p{(\linewidth - 4\tabcolsep) * \real{0.3333}}
  >{\raggedright\arraybackslash}p{(\linewidth - 4\tabcolsep) * \real{0.3333}}
  >{\raggedright\arraybackslash}p{(\linewidth - 4\tabcolsep) * \real{0.3333}}@{}}
\toprule\noalign{}
\begin{minipage}[b]{\linewidth}\raggedright
术语
\end{minipage} & \begin{minipage}[b]{\linewidth}\raggedright
英文
\end{minipage} & \begin{minipage}[b]{\linewidth}\raggedright
一句话解释
\end{minipage} \\
\midrule\noalign{}
\endhead
\bottomrule\noalign{}
\endlastfoot
智能体 & Agent & 能自己调用工具、执行多步任务来完成目标的 AI 程序 \\
智能体框架 & Agent harness &
把模型、工具、上下文、权限、会话组装起来跑循环的软件 \\
工具调用 & Tool call &
模型请求执行一个动作（读文件、跑命令、搜索）并拿到结果 \\
上下文 & Context &
模型这一次能看到的全部内容：指令、历史、文件片段、工具输出 \\
压缩 & Compaction & 上下文快满时，把过程压缩成''目标 + 结论 + 证据'' \\
权限 & Permission & 决定 Agent 能做什么、哪些动作需要你点头 \\
沙箱 & Sandbox & 限制 Agent 只能访问指定范围的机制 \\
技能 & Skill & 预先写好的领域工作流，按需加载进上下文 \\
子代理 & Subagent & 由主代理派生、专门做一件事的独立 Agent \\
计划模式 & Plan mode & 先产出计划、经你确认后再动手的工作方式 \\
会话 & Session & 一次持续的对话上下文，可保存、继续、回看 \\
工作区 & Workspace & Agent 默认能读写的项目目录 \\
\end{longtable}
}

\subsection{正文}\label{ux6b63ux6587}

\subsubsection{1.
三代编程助手：从''猜下一行''到''自己动手''}\label{ux4e09ux4ee3ux7f16ux7a0bux52a9ux624bux4eceux731cux4e0bux4e00ux884cux5230ux81eaux5df1ux52a8ux624b}

\begin{figure}
\centering
\pandocbounded{\includegraphics[keepaspectratio,alt={三代编程助手对比}]{./images/agent_vs_completion.png}}
\caption{三代编程助手对比}
\end{figure}

{\def\LTcaptype{none} % do not increment counter
\begin{longtable}[]{@{}
  >{\raggedright\arraybackslash}p{(\linewidth - 6\tabcolsep) * \real{0.2500}}
  >{\raggedright\arraybackslash}p{(\linewidth - 6\tabcolsep) * \real{0.2500}}
  >{\raggedright\arraybackslash}p{(\linewidth - 6\tabcolsep) * \real{0.2500}}
  >{\raggedright\arraybackslash}p{(\linewidth - 6\tabcolsep) * \real{0.2500}}@{}}
\toprule\noalign{}
\begin{minipage}[b]{\linewidth}\raggedright
\end{minipage} & \begin{minipage}[b]{\linewidth}\raggedright
代码补全
\end{minipage} & \begin{minipage}[b]{\linewidth}\raggedright
对话式编程
\end{minipage} & \begin{minipage}[b]{\linewidth}\raggedright
Agent 编程
\end{minipage} \\
\midrule\noalign{}
\endhead
\bottomrule\noalign{}
\endlastfoot
你给它 & 正在写的这一行 & 报错、代码片段、问题 & 目标 + 约束 +
\textbf{验收标准} \\
它给你 & 猜出的下一行 & 代码与解释 &
读文件、跑命令、改代码、再验证的\textbf{完整过程} \\
你负责 & 想清楚、跑通、验证 & 复制粘贴、串起来、验证 &
定义问题、定验收、做终审 \\
出错时 & 你改 & 你改 & 它自己看报错再改（直到你设的验收标准满足） \\
\end{longtable}
}

三者的分水岭是\textbf{闭环}：前两代的输出是''文本''，你拿到文本后还要自己去运行、观察、判断、再回来描述问题；Agent
的输出是''完成的动作''------它有眼睛（读文件、看报错）和手（改文件、跑命令），能在这个循环里自我修正。

一句话记住：\textbf{补全帮你写一行，对话帮你写一段，Agent
帮你把事情做完------但''做对''的定义仍然由你给出。}

\subsubsection{2.
harness：模型之外的四件事}\label{harnessux6a21ux578bux4e4bux5916ux7684ux56dbux4ef6ux4e8b}

\begin{figure}
\centering
\pandocbounded{\includegraphics[keepaspectratio,alt={harness 的组成}]{./images/harness_anatomy.png}}
\caption{harness 的组成}
\end{figure}

很多人以为''用了同一个模型，效果应该差不多''。实际上，同一个模型在不同
harness 里的表现可以差很远，因为决定体验的是另外四件事：

{\def\LTcaptype{none} % do not increment counter
\begin{longtable}[]{@{}
  >{\raggedright\arraybackslash}p{(\linewidth - 4\tabcolsep) * \real{0.3333}}
  >{\raggedright\arraybackslash}p{(\linewidth - 4\tabcolsep) * \real{0.3333}}
  >{\raggedright\arraybackslash}p{(\linewidth - 4\tabcolsep) * \real{0.3333}}@{}}
\toprule\noalign{}
\begin{minipage}[b]{\linewidth}\raggedright
组成
\end{minipage} & \begin{minipage}[b]{\linewidth}\raggedright
作用
\end{minipage} & \begin{minipage}[b]{\linewidth}\raggedright
少了它会怎样
\end{minipage} \\
\midrule\noalign{}
\endhead
\bottomrule\noalign{}
\endlastfoot
模型 LLM & 推理与决策 & 没有大脑，什么也做不了 \\
工具 Tools & 读写文件、跑命令、搜索 & 只能说不能做，等于聊天机器人 \\
上下文 Context & 把项目指令、历史、文件内容喂给模型 &
它不知道你的项目规范，反复问同样的问题 \\
权限 Permission & 决定能不能写、要不要审批 &
要么什么都不敢做，要么风险失控 \\
会话 Session & 记录全部动作，可继续、可回看 &
每次从零开始，出了问题无法追责 \\
\end{longtable}
}

\begin{quote}
类比：模型是发动机，harness
是底盘、变速箱、刹车和仪表盘。发动机再好，没有刹车也不敢上路。
\end{quote}

\subsubsection{3. 一次 Agent
循环的解剖}\label{ux4e00ux6b21-agent-ux5faaux73afux7684ux89e3ux5256}

\begin{figure}
\centering
\pandocbounded{\includegraphics[keepaspectratio,alt={一次 Agent 循环}]{./images/agent_loop.png}}
\caption{一次 Agent 循环}
\end{figure}

循环只有五个动作，但每一步都有讲究：

\begin{enumerate}
\def\labelenumi{\arabic{enumi}.}
\tightlist
\item
  \textbf{读}：看目录、读文件、看报错。它看到什么，取决于你给的工作区与上下文。
\item
  \textbf{想}：拆解任务、定计划。复杂任务建议先进入\textbf{计划模式}，让你确认后再动手。
\item
  \textbf{做}：调用工具------改代码、跑脚本、装依赖、查文档。
\item
  \textbf{看}：读回执行结果。这一步是 Agent
  与''聊天''最大的区别：它必须面对真实输出。
\item
  \textbf{判}：达标了吗？达标的判据\textbf{不是''代码跑起来了''，而是你事先给的验收标准}。
\end{enumerate}

用一个具体任务走一遍。假设你有一份成绩矩阵，想算''每个学生的平均分''，但代码写成了按列求平均：

\begin{itemize}
\tightlist
\item
  读：Agent 打开脚本，看到 scores 的形状是 (4, 3)、代码里写着
  mean(axis=0)。
\item
  想：它需要判断你的\textbf{意图}------``每个学生''对应行还是列。
\item
  做：它写出最小复现，分别打印两种 axis 的结果和形状。
\item
  看：形状一个 (3,) 一个 (4,)，它意识到意图与实现不一致。
\item
  判：按你给的验收标准（形状必须是
  (4,)、数值与手算一致）给出结论与修复。
\end{itemize}

\begin{quote}
关键认识：\textbf{循环的终止条件由你定义。} 没有验收标准，Agent
就会以''看起来能跑''作为终点。
\end{quote}

\subsubsection{4. 上下文：Agent
的短期记忆}\label{ux4e0aux4e0bux6587agent-ux7684ux77edux671fux8bb0ux5fc6}

\begin{figure}
\centering
\pandocbounded{\includegraphics[keepaspectratio,alt={上下文窗口与压缩}]{./images/context_window.png}}
\caption{上下文窗口与压缩}
\end{figure}

上下文窗口再大也会满。一次长任务里，系统提示、对话历史、读入的文件、工具输出会迅速堆积；快满时
harness
会做\textbf{压缩（compaction）}，把过程浓缩成目标、结论与关键证据。

三条实用建议：

\begin{enumerate}
\def\labelenumi{\arabic{enumi}.}
\tightlist
\item
  \textbf{给最小可复现}：一段 20 行的复现脚本，比 2000
  行项目源码更有用。
\item
  \textbf{让它自己读文件}：说''读 src/clean.py 里 load\_data
  函数''比把代码整段贴进对话更省上下文，也更准确。
\item
  \textbf{长任务主动收口}：阶段结束时让它总结''已确认的结论 +
  未解决的问题 + 下一步''，这既是压缩，也是你检查方向的时机。
\end{enumerate}

\paragraph{案例卡 1：一个''看起来都对''的错误------axis
用反了}\label{ux6848ux4f8bux5361-1ux4e00ux4e2aux770bux8d77ux6765ux90fdux5bf9ux7684ux9519ux8befaxis-ux7528ux53cdux4e86}

\textbf{目标}：认识''静默数值错误''------代码能跑、数字看着合理，但结论是错的。

\textbf{代码}

\begin{Shaded}
\begin{Highlighting}[]
\ImportTok{import}\NormalTok{ numpy }\ImportTok{as}\NormalTok{ np}
\NormalTok{np.set\_printoptions(precision}\OperatorTok{=}\DecValTok{2}\NormalTok{, suppress}\OperatorTok{=}\VariableTok{True}\NormalTok{)}

\NormalTok{scores }\OperatorTok{=}\NormalTok{ np.array([[}\DecValTok{82}\NormalTok{, }\DecValTok{90}\NormalTok{, }\DecValTok{78}\NormalTok{],}
\NormalTok{                   [}\DecValTok{60}\NormalTok{, }\DecValTok{65}\NormalTok{, }\DecValTok{70}\NormalTok{],}
\NormalTok{                   [}\DecValTok{95}\NormalTok{, }\DecValTok{88}\NormalTok{, }\DecValTok{92}\NormalTok{],}
\NormalTok{                   [}\DecValTok{70}\NormalTok{, }\DecValTok{72}\NormalTok{, }\DecValTok{68}\NormalTok{]], dtype}\OperatorTok{=}\BuiltInTok{float}\NormalTok{)}

\NormalTok{by\_exam }\OperatorTok{=}\NormalTok{ scores.mean(axis}\OperatorTok{=}\DecValTok{0}\NormalTok{)}
\NormalTok{by\_student }\OperatorTok{=}\NormalTok{ scores.mean(axis}\OperatorTok{=}\DecValTok{1}\NormalTok{)}

\BuiltInTok{print}\NormalTok{(}\StringTok{"by\_exam   :"}\NormalTok{, by\_exam, }\StringTok{"shape"}\NormalTok{, by\_exam.shape)}
\BuiltInTok{print}\NormalTok{(}\StringTok{"by\_student:"}\NormalTok{, by\_student, }\StringTok{"shape"}\NormalTok{, by\_student.shape)}
\BuiltInTok{print}\NormalTok{(}\StringTok{"mean of all:"}\NormalTok{, }\BuiltInTok{round}\NormalTok{(}\BuiltInTok{float}\NormalTok{(scores.mean()), }\DecValTok{3}\NormalTok{))}
\BuiltInTok{print}\NormalTok{(}\StringTok{"sum check :"}\NormalTok{, }\BuiltInTok{round}\NormalTok{(}\BuiltInTok{float}\NormalTok{(by\_exam.}\BuiltInTok{sum}\NormalTok{()), }\DecValTok{3}\NormalTok{), }\BuiltInTok{round}\NormalTok{(}\BuiltInTok{float}\NormalTok{(by\_student.}\BuiltInTok{sum}\NormalTok{()), }\DecValTok{3}\NormalTok{))}
\end{Highlighting}
\end{Shaded}

\textbf{运行结果（已运行核验）}

\begin{Shaded}
\begin{Highlighting}[]
\NormalTok{by\_exam   : [76.75 78.75 77.  ] shape (3,)}
\NormalTok{by\_student: [83.33 65.   91.67 70.  ] shape (4,)}
\NormalTok{mean of all: 77.5}
\NormalTok{sum check : 232.5 310.0}
\end{Highlighting}
\end{Shaded}

\textbf{讲解}

\begin{itemize}
\tightlist
\item
  两个结果都是''合理的分数''，粗看都像平均分------这正是静默错误可怕的地方；
\item
  形状是第一个判据：想算''每个学生''就该得到 4 个数，得到 3
  个说明方向错了；
\item
  总和是第二个判据：by\_exam 的均值之和 232.5 只是总分 930
  的三分之一，by\_student 的 310.0 才是''每人平均分之和''；
\item
  只打印数字、不做判据检查，这类错误可以躲过一整天。
\end{itemize}

\textbf{主要用法 / API}

{\def\LTcaptype{none} % do not increment counter
\begin{longtable}[]{@{}ll@{}}
\toprule\noalign{}
用法 & 说明 \\
\midrule\noalign{}
\endhead
\bottomrule\noalign{}
\endlastfoot
array.mean(axis=k) & 沿第 k 维求平均：k=0 消掉行、k=1 消掉列 \\
array.shape & 判断结果方向是否符合意图的第一判据 \\
np.set\_printoptions & 控制打印精度，避免误以为''数不一样'' \\
\end{longtable}
}

\textbf{常见错误}

\begin{enumerate}
\def\labelenumi{\arabic{enumi}.}
\tightlist
\item
  只把结果打印出来看一眼''像不像''，不用形状与守恒量检查；
\item
  把 axis=0 记成''对每一行''（其实是沿行方向压缩，得到每列的值）；
\item
  让 Agent 修 bug
  时只说''结果不对''，不说\textbf{期望的形状与判据}，它只能猜。
\end{enumerate}

\textbf{拓展}

\begin{itemize}
\tightlist
\item
  把 scores 换成 (3, 4, 5) 的三维数组，先说出 mean(axis=1)
  的形状，再运行验证；
\item
  让 Agent 写一个 3 行的断言函数，把''形状必须是 (4,)``固化成测试；
\item
  思考：如果这份数据是某次实验的唯一结果，这个错误会怎样影响结论？
\end{itemize}

\subsubsection{5.
人机分工：什么交给它，什么必须自己管}\label{ux4ebaux673aux5206ux5de5ux4ec0ux4e48ux4ea4ux7ed9ux5b83ux4ec0ux4e48ux5fc5ux987bux81eaux5df1ux7ba1}

{\def\LTcaptype{none} % do not increment counter
\begin{longtable}[]{@{}lll@{}}
\toprule\noalign{}
任务类型 & 建议 & 原因 \\
\midrule\noalign{}
\endhead
\bottomrule\noalign{}
\endlastfoot
样板代码、格式转换、批量重命名 & 交给 Agent & 规则明确、易验证 \\
写测试、补边界条件、重构 & 交给 Agent（你审 diff） & 有客观判据 \\
查文档、读陌生代码库、写注释 & 交给 Agent & 信息检索型工作 \\
问题定义与建模 & \textbf{人主导} & 决定''做什么''，Agent 无法替你决定 \\
验收标准与数据判据 & \textbf{人主导} & 决定''什么叫做对'' \\
结论与解释 & \textbf{人主导}（Agent 只能起草） & 涉及因果、局限与责任 \\
\end{longtable}
}

三条铁律：\textbf{问题定义归人、验收标准归人、结论归人。} Agent
可以在这三件事上给建议，但不能替你签字。

\subsubsection{6. 权限与安全：让 Agent
在笼子里干活}\label{ux6743ux9650ux4e0eux5b89ux5168ux8ba9-agent-ux5728ux7b3cux5b50ux91ccux5e72ux6d3b}

\begin{figure}
\centering
\pandocbounded{\includegraphics[keepaspectratio,alt={权限档位与审批}]{./images/permission_flow.png}}
\caption{权限档位与审批}
\end{figure}

主流 harness 都提供分档权限，常见三档：

{\def\LTcaptype{none} % do not increment counter
\begin{longtable}[]{@{}
  >{\raggedright\arraybackslash}p{(\linewidth - 4\tabcolsep) * \real{0.3333}}
  >{\raggedright\arraybackslash}p{(\linewidth - 4\tabcolsep) * \real{0.3333}}
  >{\raggedright\arraybackslash}p{(\linewidth - 4\tabcolsep) * \real{0.3333}}@{}}
\toprule\noalign{}
\begin{minipage}[b]{\linewidth}\raggedright
档位
\end{minipage} & \begin{minipage}[b]{\linewidth}\raggedright
含义
\end{minipage} & \begin{minipage}[b]{\linewidth}\raggedright
适用场景
\end{minipage} \\
\midrule\noalign{}
\endhead
\bottomrule\noalign{}
\endlastfoot
read-only & 只能读，不能改 & 第一次接触陌生项目、只想让它解释代码 \\
workspace-write & 只能改工作区内的文件 &
日常开发、课程作业（\textbf{推荐默认}） \\
danger-full-access & 全盘可写、命令直通 & 受控环境下的高级用法，慎用 \\
\end{longtable}
}

配套机制是\textbf{审批}：涉及写文件、跑命令、装依赖时，界面会先问你；你的每一次允许或拒绝都会写进会话记录，可回看、可追责。

数据红线（请直接记住）：

\begin{itemize}
\tightlist
\item
  不要把隐私数据（身份证、手机号、成绩单、医疗记录）贴进对话；
\item
  不要把涉密或未发表的数据交给云端模型；
\item
  不要把别人的 API Key、密码、令牌放进提示词或代码；
\item
  处理真实数据前，先想清楚''这份数据能不能离开这台机器''。
\end{itemize}

\subsubsection{7.
四种典型失败模式与对策}\label{ux56dbux79cdux5178ux578bux5931ux8d25ux6a21ux5f0fux4e0eux5bf9ux7b56}

{\def\LTcaptype{none} % do not increment counter
\begin{longtable}[]{@{}
  >{\raggedright\arraybackslash}p{(\linewidth - 4\tabcolsep) * \real{0.3333}}
  >{\raggedright\arraybackslash}p{(\linewidth - 4\tabcolsep) * \real{0.3333}}
  >{\raggedright\arraybackslash}p{(\linewidth - 4\tabcolsep) * \real{0.3333}}@{}}
\toprule\noalign{}
\begin{minipage}[b]{\linewidth}\raggedright
失败模式
\end{minipage} & \begin{minipage}[b]{\linewidth}\raggedright
长什么样
\end{minipage} & \begin{minipage}[b]{\linewidth}\raggedright
对策
\end{minipage} \\
\midrule\noalign{}
\endhead
\bottomrule\noalign{}
\endlastfoot
幻觉 API & 编出不存在的函数、参数、配置项 &
要求它运行验证；自己查一眼官方文档 \\
静默数值错误 & 能跑、数字合理，其实方向或单位错了 &
事先给判据：形状、量纲、守恒量、解析解对照 \\
过度自信的结论 & 把相关说成因果、不给不确定性 &
要求给证据、给区间、给反例，结论自己写 \\
改坏别处 & 顺手重构、动了不该动的文件 & 小步提交、看清
diff、跑测试再接受 \\
\end{longtable}
}

\begin{quote}
这四种失败都不是''模型不够聪明''，而是\textbf{缺验证}。第 03
节会给出一个固定的验收工作流。
\end{quote}

\subsubsection{8. 上手体验（5
分钟）}\label{ux4e0aux624bux4f53ux9a8c5-ux5206ux949f}

如果你已经按第 02 节装好了 DeepSeek
Harness，现在就可以试一次。任务：让它解释本节的 axis 案例。

\begin{Shaded}
\begin{Highlighting}[]
\NormalTok{我在做课程第 1 章的练习题，这段代码想算"每个学生的平均分"，但我不确定对不对：}

\NormalTok{scores = np.array([[82, 90, 78], [60, 65, 70], [95, 88, 92], [70, 72, 68]], dtype=float)}
\NormalTok{by\_exam = scores.mean(axis=0)}

\NormalTok{请你：}
\NormalTok{1) 指出这段代码算出来的到底是什么（给出形状）；}
\NormalTok{2) 写出正确写法，并给出一个最小复现脚本；}
\NormalTok{3) 加一条断言，保证结果形状符合"每个学生"的意图。}
\NormalTok{不要修改我工作区里的其他文件。}
\end{Highlighting}
\end{Shaded}

观察三个点：

\begin{enumerate}
\def\labelenumi{\arabic{enumi}.}
\tightlist
\item
  它有没有\textbf{先读文件、先运行}，还是直接给答案；
\item
  它给出的断言是不是可执行、可验证；
\item
  它有没有解释''为什么 axis=0 得到的是每列''，而不是只给结论。
\end{enumerate}

\subsection{本节小结}\label{ux672cux8282ux5c0fux7ed3}

\begin{enumerate}
\def\labelenumi{\arabic{enumi}.}
\tightlist
\item
  三代工具的分水岭是\textbf{闭环}：能否自己执行、观察、修正。
\item
  harness = 模型 + 工具 + 上下文 + 权限 + 会话，后四件决定体验。
\item
  循环五步：读、想、做、看、判；\textbf{终止条件由你的验收标准定义}。
\item
  上下文是最稀缺资源：给最小复现、让它自己读文件、长任务主动收口。
\item
  权限三档 + 审批 + 数据红线；问题定义、验收标准、结论永远归人。
\end{enumerate}

\subsection{思考题}\label{ux601dux8003ux9898}

\begin{enumerate}
\def\labelenumi{\arabic{enumi}.}
\tightlist
\item
  同一段代码，让''对话式助手''和''Agent''分别修一个
  bug，你能观察到哪些行为差异？
\item
  为什么''代码能跑通''不能作为验收标准？请举一个本课程里的例子。
\item
  如果上下文快满了，你会保留哪三类信息？
\item
  什么情况下应该主动把权限降到 read-only？
\item
  同学说''Agent
  写的报告结论和我们实验一致，所以没问题''。请指出这句话的逻辑漏洞。
\end{enumerate}

\subsection{练习与延伸阅读}\label{ux7ec3ux4e60ux4e0eux5ef6ux4f38ux9605ux8bfb}

\begin{itemize}
\tightlist
\item
  练习：第 11 章 \texttt{exercises/} 第 1--4 题（概念题）；
\item
  延伸：DeepSeek Harness 官方仓库与文档（见 \texttt{references.md}）；
\item
  下一节：\texttt{02-安装与配置.md}（三条安装路径 + API Key
  申请与付费）。
\end{itemize}

\end{document}
